\section{Feilhåndtering og logging}
\label{sec:errors}

I et mellomvare-API som videresender kall til to eller flere ulike eksterne videomotorer kan feil oppstå i flere lag. Klienten kan sende ugyldige forespørsler, backend-tjenester kan returnere uventede statuskoder eller være utilgjengelige, og konfigurasjonsverdier kan mangle. For å gjøre slike feil håndterbare for både klientutviklere og driftspersonell ble det innført en sentralisert strategi der domenespesifikke exceptions defineres ett sted, controllers fanger ikke feil selv, og alle uhåndterte exceptions blir oversatt til et felles svarformat av en global handler.

\subsection{Strategi: én sentral handler}
\label{subsec:errors-strategi}

Tilnærmingen bygger på tre prinsipper. For det første skal controllerne ikke fange exceptions; feil får i stedet propagere til middleware-laget. Dette holder controller-koden lesbar og hindrer at hvert endepunkt må implementere samme \texttt{try/catch}-logikk på nytt. For det andre skal all feilrespons ta utgangspunkt i RFC~7807, \emph{Problem Details for HTTP APIs} \cite{rfc7807}, slik at klienter får presis og forståelig informasjon om hva som gikk galt. For det tredje skal interne detaljer ikke eksponeres. Feilmeldinger med statuskode \texttt{5xx} erstattes med generisk tekst, mens detaljert informasjon logges internt (jf. kravspesifikasjonen \cite{hdo_kravspec}).

ASP.NET Core 10 har innebygd støtte for dette mønsteret via \texttt{IExceptionHandler}. I oppstartsekvensen registreres handleren som en tjeneste, og middleware-kjeden aktiverer den rett etter routing:

\begin{lstlisting}[language={[Sharp]C}, caption={Registrering av den globale feil-handleren.}, label={lst:register-handler}]
// HDO-VideoAPI/Extensions/VideoApiCoreExtensions.cs

services.AddProblemDetails();
services.AddExceptionHandler<GlobalExceptionHandler>();

// HDO-VideoAPI/Extensions/ApplicationExtensions.cs

app.UseExceptionHandler();
\end{lstlisting}

\subsection{Domenespesifikke exceptions}
\label{subsec:errors-exceptions}

For å unngå at provider-koden kaster generiske \texttt{Exception}-objekter som ville endt opp som \texttt{500 Internal Server Error}, defineres fire egendefinerte exceptions i \texttt{Exceptions/VideoApiExceptions.cs}. Hver type representerer en kategori av feil som tilsvarer en bestemt HTTP-status, og bærer med seg strukturert kontekst (ressurs-ID, tjenestenavn, ekstern statuskode, konfigurasjonsnøkkel) som senere kan eksponeres til klienten eller loggføres.

Løsningen er modulær og kan enkelt utvides med nye domenespesifikke exceptions ved behov, uten endringer i eksisterende håndteringslogikk. Tabell~\ref{tab:domain-exceptions} oppsummerer typene.

\begin{table}[htbp]
\centering
\caption{Domenespesifikke exceptions og tilhørende HTTP-statuskoder.}
\label{tab:domain-exceptions}
\small
\begin{tabular}{l l l}
\toprule
\textbf{Exception} & \textbf{HTTP-status} & \textbf{Brukes når \dots} \\
\midrule
\texttt{BadRequestException}      & \texttt{400 Bad Request}          & klienten sender ugyldige verdier \\
\texttt{NotFoundException}        & \texttt{404 Not Found}            & rom eller en booking finnes ikke \\
\texttt{ExternalServiceException} & \texttt{502 Bad Gateway}          & AntMedia eller NHN feiler \\
\texttt{ConfigurationException}   & \texttt{503 Service Unavailable}  & ugyldig/manglende konfigurasjon \\
\bottomrule
\end{tabular}
\end{table}

\newpage
% I praksis kastes unntakene fra provider-laget og service-laget der feilen faktisk oppdages. Kodeliste~\ref{lst:notfound-throw} viser hvordan AntMedia-provideren oversetter en HTTP-status \texttt{404 Not Found} fra ekstern tjeneste til \texttt{NotFoundException}, og hvordan andre feil pakkes inn som \texttt{ExternalServiceException} med den opprinnelige statuskoden bevart:

I praksis kastes exceptions fra provider-laget der feilen faktisk oppdages. \texttt{AntMediaProvider} benytter en privat hjelpemetode, \texttt{EnsureSuccessOrThrowAsync}, som kalles etter hvert kall mot den eksterne tjenesten og oversetter HTTP-feil til domenespesifikke exceptions: en \texttt{404 Not Found} blir til \texttt{NotFoundException}, mens andre feil pakkes inn som \texttt{ExternalServiceException} med den opprinnelige statuskoden bevart (kodeliste~\ref{lst:notfound-throw}).

\begin{lstlisting}[language={[Sharp]C}, caption={Kasting av domenespesifik exception i \texttt{AntMediaProvider}.}, label={lst:notfound-throw}]
// HDO-VideoAPI/Providers/AntMedia/AntMediaProvider.cs

private static async Task EnsureSuccessOrThrowAsync( ... )
{
    // Vellykket svar - ingen feil å håndtere
    if (response.IsSuccessStatusCode)
        return;

    // 404 oversettes til NotFoundException
    if (response.StatusCode == HttpStatusCode.NotFound && roomId is not null)
        throw new NotFoundException("Room", roomId);

    // Alle andre feil - bevarer opprinnelig statuskode og responstekst
    var status = (int)response.StatusCode;
    var body = await response.Content.ReadAsStringAsync(ct);
    throw new ExternalServiceException(
        "AntMedia",
        $"AntMedia API returned {status} for {operation}. {body}",
        status);
}
\end{lstlisting}

\subsection{Den globale handleren}
\label{subsec:errors-handler}

% \texttt{GlobalExceptionHandler} implementerer \texttt{IExceptionHandler.TryHandleAsync}. Metoden mottar det fangede exception-et, tilordner det en HTTP-statuskode ved hjelp av et \texttt{switch}-uttrykk (se kodeliste \ref{lst:mapping}), registrerer en metrikk, logger hendelsen med korrekt alvorlighetsgrad og returnerer et \texttt{ProblemDetails}-objekt som JSON til klienten. Koblingen mellom exceptions og HTTP-statuskoder er sentralisert, slik at en ny exception-type kun krever én kodeendring:

\texttt{GlobalExceptionHandler} implementerer \texttt{IExceptionHandler.TryHandleAsync}. Metoden mottar det kastede exception-et, slår opp riktig HTTP-statuskode via et \texttt{switch}-uttrykk (kodeliste~\ref{lst:mapping}), registrerer en metrikk, logger hendelsen med korrekt alvorlighetsgrad og returnerer et \texttt{ProblemDetails}-objekt som JSON til klienten. Koblingen mellom exceptions og statuskoder er sentralisert, slik at en ny exception-type kun krever én kodeendring.

\begin{lstlisting}[language={[Sharp]C}, caption={Tilordning fra exception-type til HTTP-statuskode i \texttt{GlobalExceptionHandler}}, label={lst:mapping}]
// HDO-VideoAPI/Middleware/GlobalExceptionHandler.cs

public async ValueTask<bool> TryHandleAsync( ... )
{
    var (statusCode, errorResponse) = exception switch
    {
        BadRequestException        => (400, new { error = "Bad Request" }),
        ArgumentException          => (400, new { error = "Invalid Argument" }),
        NotFoundException          => (404, new { error = "Resource Not Found" }),
        OperationCanceledException => (499, new { error = "Request was cancelled" }),
        NotSupportedException      => (501, new { error = "Not Supported" }),
        ExternalServiceException   => (502, new { error = "External Service Error" }),
        ConfigurationException     => (503, new { error = "Service Unavailable due to configuration error" }),
        _                          => (500, new { error = "An unexpected error occurred." })
    };

    // ... (metrikk, logging, RFC 7807 ProblemDetails-konstruksjon og JSON-respons)
}
\end{lstlisting}

I tillegg til de fire domenespesifikke exception-typene håndteres også \texttt{ArgumentException}, \texttt{OperationCanceledException} og \texttt{NotSupportedException} eksplisitt, slik at vanlige rammeverks-exceptions ikke eksponeres som \texttt{500}-svar. Den avsluttende \texttt{\_}-armen fungerer som et sikkerhetsnett og returnerer \texttt{500 Internal Server Error} med en generisk feilmelding for alle exceptions som ikke fanges av de øvrige armene.

\subsection{Svarformat: Problem Details}
\label{subsec:errors-problemdetails}

Selve responsen bygges som et \texttt{ProblemDetails}-objekt og serialiseres med \texttt{Content-Type: application/problem+json}, slik RFC~7807 krever \cite{rfc7807}. For \texttt{5xx}-feil overskrives \texttt{Detail}-feltet med en generisk tekst for å unngå at sensitive detaljer eksponeres mot klienten, mens \texttt{traceId} alltid legges inn i utvidelsene slik at en loggoppføring kan kobles til et bestemt feilsvar i ettertid.

\begin{lstlisting}[language={[Sharp]C}, caption={Bygging av \texttt{ProblemDetails}-respons.}, label={lst:problemdetails}]
// HDO-VideoAPI/Middleware/GlobalExceptionHandler.cs

public async ValueTask<bool> TryHandleAsync( ... )
{
    // ... (statuskode og feilmelding fra exception-type, metrikk og logging)

    // Bygge ProblemDetails objekt
    var problemDetails = new ProblemDetails
    {
        Status   = statusCode,
        Title    = errorResponse.error,
        Type     = $"https://httpstatuses.com/{statusCode}",
        Detail   = statusCode >= 500
                     ? "An internal error occurred. See the logs for more detail."
                     : exception.Message,
        Instance = httpContext.Request.Path
    };
    problemDetails.Extensions["traceId"] = httpContext.TraceIdentifier;

    // Bevarer ekstern statuskode for å skille mellom ulike 502-årsaker
    if (exception is ExternalServiceException ext && ext.StatusCode.HasValue)
        problemDetails.Extensions["externalStatusCode"] = ext.StatusCode.Value;
    
    httpContext.Response.StatusCode = statusCode;
    await httpContext.Response.WriteAsJsonAsync(
        problemDetails, options: null,
        contentType: "application/problem+json", cancellationToken);
}
\end{lstlisting}

For \texttt{ExternalServiceException} legges også feltet \texttt{externalStatusCode} til, slik at klienten kan skille mellom ulike årsaker til et \texttt{502 Bad Gateway}-svar --- for eksempel om AntMedia returnerte \texttt{401 Unauthorized} eller NHN ga \texttt{500 Internal Server Error}. Et typisk svar fra et oppslag på et ikke-eksisterende rom ser slik ut:

\begin{lstlisting}[caption={Eksempel på Problem Details-respons fra API-et (HTTP 404).}, label={lst:problem-example}]
HTTP/1.1 404 Not Found
Content-Type: application/problem+json

{
  "status": 404,
  "title":  "Resource Not Found",
  "type":   "https://httpstatuses.com/404",
  "detail": "Room with ID 'abc123' was not found.",
  "instance": "/api/rooms/antmedia:abc123",
  "traceId":  "00-7c2f...-01"
}
\end{lstlisting}

\subsection{Logging og metrikker}
\label{subsec:errors-observability}

Hver fangede exception registreres som et inkrement i Prometheus-telleren \texttt{hdo\_errors\_total}, etiketten settes til \texttt{exception.GetType().Name}. Det gir et grunnlag for å se hvilke typer feil som dominerer i produksjon uten å måtte parse logger. Logging skjer med to alvorlighetsgrader avhengig av statuskoden, slik at \texttt{4xx}-feil ikke fyller opp varsel-kanalen for ekte produksjonsfeil.

\begin{lstlisting}[language={[Sharp]C}, caption={Differensiert logging i den globale håndtereren.}, label={lst:logging}]
metrics.ErrorOccurred(exception.GetType().Name);

if (statusCode >= 500)
    logger.LogError(exception, "An unhandled exception occurred.");
else
    logger.LogWarning(exception, "A handled exception occurred");
\end{lstlisting}

I tråd med krav O1.1 \cite{hdo_kravspec} logges aldri API-nøkler, JWT-tokens eller annen hemmelig informasjon. Exception-typene er bevisst utformet slik at meldingsfeltet kun inneholder ressurs-ID-er og tjenestenavn, og ikke potensielt sensitive verdier.

\subsection{Klientsidens behandling av Problem Details}
\label{subsec:errors-client}

Demo-frontenden speiler kontrakten i en liten hjelper i \texttt{Frontend/src/utils/errors.ts}. Funksjonen \texttt{apiCall} pakker hver \texttt{fetch}-respons og kaster en \texttt{ApiError} som bærer hele \texttt{ProblemDetails}-objektet. Det betyr at sidekomponentene kan vise en brukervennlig tekst med både feilbeskrivelse og \texttt{traceId} uten å bygge parsing-logikk selv:

\begin{lstlisting}[language=JavaScript, caption={Klientsidens tolkning av Problem Details (forenklet).}, label={lst:client-error}]
// Frontend/src/utils/errors.ts
export async function apiCall<T>(fn: () => Promise<Response>): Promise<T> {
  const res = await fn();
  if (!res.ok) {
    const problem = await res.json().catch(() => ({
      status: res.status, title: res.statusText, detail: '', instance: '', traceId: ''
    }));
    throw new ApiError(problem);
  }
  // ...
}

export function getErrorMessage(err: unknown): string {
  if (err instanceof ApiError) {
    const trace = err.problem.traceId ? ` (Trace: ${err.problem.traceId})` : '';
    return `${err.problem.detail || err.problem.title}${trace}`;
  }
  return err instanceof Error ? err.message : 'An unexpected error occurred.';
}
\end{lstlisting}

Det fullstendige bildet er en feilhåndteringskjede der provideren kaster en domenespesifikk exception, middleware-komponenten oversetter denne til et standardisert \texttt{problem+json}-svar, en teller inkrementeres, og klienten viser en menneskelig lesbar melding med referanse til loggingen. Modellen gir tre konkrete fordeler: Nye exception-typer kan legges til uten å endre verken controller- eller provider-kode, klienter trenger kun én parser for alle feiltyper, og driftspersonell kan korrelere en synlig feil i brukergrensesnittet med en spesifikk loggpost via \texttt{traceId}.
